% =========================================================
% Chapter 3 — Stationary time series
% Author : Aymen Ben Brik (aymen.benbrik@esprit.tn)
% =========================================================

\chapter{Stationary time series}
\label{ch:stationarity}

\section*{Introduction}
\addcontentsline{toc}{section}{Introduction}

After decomposing the observed series into trend, seasonality and
residual (Chapter~\ref{ch:decomposition}), we usually want the residual
to be \emph{stationary} so that the modelling tools of
Chapter~\ref{ch:arma} apply. This chapter formalises stationarity,
introduces the autocovariance and autocorrelation functions (ACF and
PACF), states the Wold decomposition theorem, and provides
\emph{white-noise tests} (Box--Pierce, Ljung--Box) to validate model
residuals.

% =========================================================
\section{Stationarity}
\label{sec:stationarity}

\subsection{Strict and weak stationarity}

\begin{defbox}[Strict (strong) stationarity]
\label{def:strict-stat}
A process $\{X_t\}_{t \in \Z}$ is \emph{strictly stationary} if, for
every integer $k$, every $t_1 < \dots < t_k$ and every $h \in \Z$,
\[
  (X_{t_1}, X_{t_2}, \dots, X_{t_k})
  \overset{\text{d}}{=}
  (X_{t_1 + h}, X_{t_2 + h}, \dots, X_{t_k + h}).
\]
\end{defbox}

\begin{defbox}[Weak (covariance) stationarity]
\label{def:weak-stat}
A process $\{X_t\}$ with $\E[X_t^2] < \infty$ is \emph{weakly
stationary} if there exist $\mu \in \R$ and a function
$\gamma : \Z \to \R$ such that
\[
  \E[X_t] = \mu, \qquad
  \Cov(X_t, X_{t+h}) = \gamma(h) \quad \text{for all } t, h \in \Z.
\]
The function $\gamma$ is the \emph{autocovariance function} (ACVF).
\end{defbox}

\begin{propbox}[Strict $\Rightarrow$ weak (under $L^2$)]
\label{prop:strict-weak}
If $\{X_t\}$ is strictly stationary with $\E[X_t^2] < \infty$, then
$\{X_t\}$ is weakly stationary.
\end{propbox}

\begin{proof}
By strict stationarity, $X_t \overset{\text{d}}{=} X_0$, so
$\E[X_t] = \E[X_0] = \mu$ does not depend on $t$. Similarly,
$(X_t, X_{t+h}) \overset{\text{d}}{=} (X_0, X_h)$, hence
$\Cov(X_t, X_{t+h}) = \Cov(X_0, X_h)$ depends only on $h$.
\end{proof}

\begin{rembox}[The converse is false in general]
A process can be weakly stationary without being strictly stationary
if it has time-varying higher-order moments. For Gaussian processes
(only mean and covariance matter), the two notions coincide.
\end{rembox}

\subsection{Two canonical examples}

\begin{defbox}[White noise]
\label{def:wn}
$(\varepsilon_t)$ is a \emph{white noise} $\WN(0, \sigma^2)$ if
$\E[\varepsilon_t] = 0$, $\Var[\varepsilon_t] = \sigma^2$ and
$\Cov(\varepsilon_t, \varepsilon_{t+h}) = 0$ for $h \neq 0$.
A \emph{strong white noise} requires the $\varepsilon_t$ to be
i.i.d.; a Gaussian white noise has $\varepsilon_t \iid \mathcal{N}(0, \sigma^2)$.
\end{defbox}

\begin{exbox}[Random walk is \emph{not} stationary]
\label{ex:random-walk}
Let $\varepsilon_t \iid \WN(0, \sigma^2)$ and
$X_t = X_{t-1} + \varepsilon_t$ with $X_0 = 0$. Then
$X_t = \sum_{s = 1}^{t} \varepsilon_s$ and
\[
  \E[X_t] = 0,
  \quad
  \Var[X_t] = t\, \sigma^2,
  \quad
  \Cov(X_t, X_{t+h}) = t\, \sigma^2 \text{ for } h \geq 0.
\]
The variance grows linearly in $t$, so $X_t$ is non-stationary. This
is the canonical example used by unit-root tests in
Chapter~\ref{ch:nonstationarity}.
\end{exbox}

% =========================================================
\section{Autocovariance and autocorrelation}
\label{sec:acf}

\begin{defbox}[ACVF and ACF]
For a weakly stationary process with mean $\mu$:
\[
  \gamma(h) = \Cov(X_t, X_{t+h}),
  \qquad
  \rho(h) = \frac{\gamma(h)}{\gamma(0)} \;\in\; [-1, 1].
\]
\end{defbox}

\begin{propbox}[Properties of $\gamma$]
\label{prop:acvf}
\begin{enumerate}
  \item $\gamma(0) = \Var[X_t] \geq 0$;
  \item $\abs{\gamma(h)} \leq \gamma(0)$ (Cauchy--Schwarz);
  \item $\gamma$ is even: $\gamma(-h) = \gamma(h)$;
  \item $\gamma$ is positive semidefinite: for any $n$ and any
        $a_1, \dots, a_n \in \R$,
        $\sum_{i, j = 1}^{n} a_i\, a_j\, \gamma(t_i - t_j) \geq 0$.
\end{enumerate}
\end{propbox}

\begin{proof}
(1) and (3) follow from the definition. (2) follows from Cauchy--Schwarz
applied to $X_t - \mu$ and $X_{t+h} - \mu$. (4) is the standard fact
that any matrix of the form $\Sigma_{ij} = \gamma(t_i - t_j)$ is the
covariance matrix of $(X_{t_1}, \dots, X_{t_n})$, hence positive
semidefinite.
\end{proof}

\subsection{Sample ACF and Bartlett bands}

\begin{defbox}[Sample ACVF and ACF]
For $X_1, \dots, X_T$ with sample mean $\bar X$:
\[
  \widehat\gamma(h) = \frac{1}{T}\sum_{t = 1}^{T - h}
    (X_t - \bar X)(X_{t+h} - \bar X),
  \qquad
  \widehat\rho(h) = \frac{\widehat\gamma(h)}{\widehat\gamma(0)}.
\]
\end{defbox}

\begin{thbox}[Bartlett's formula for white noise]
\label{th:bartlett}
If $\{X_t\}$ is a white noise with finite fourth moment, then for any
fixed $h \geq 1$:
\[
  \sqrt{T}\,\widehat\rho(h) \;\xrightarrow{d}\; \mathcal{N}(0, 1).
\]
The standard test bands $\pm 1.96 / \sqrt{T}$ on the ACF plot are
asymptotic 95\% confidence intervals \emph{under the white-noise null}.
\end{thbox}

\begin{rembox}
Bartlett's formula gives a more general result for non-white processes
(involves a sum over the squared autocorrelations); see Brockwell \&
Davis, \emph{Time Series: Theory and Methods}, Chapter 7.
\end{rembox}

% =========================================================
\section{Partial autocorrelation}
\label{sec:pacf}

\begin{defbox}[Partial autocorrelation function]
\label{def:pacf}
The PACF $\alpha(h)$ at lag $h$ is the correlation between $X_t$ and
$X_{t+h}$ \emph{after removing the linear effect} of the intermediate
observations $X_{t+1}, \dots, X_{t+h-1}$:
\[
  \alpha(h) = \Corr\bigl(X_t - \widehat X_t,\; X_{t+h} - \widehat X_{t+h}\bigr),
\]
where $\widehat X_t$ and $\widehat X_{t+h}$ are the linear projections
of $X_t$ and $X_{t+h}$ on the span of $\{X_{t+1}, \dots, X_{t+h-1}\}$.
\end{defbox}

\begin{rembox}[Recursive computation]
The PACF is computed by the \emph{Durbin--Levinson} algorithm: it
solves the Yule--Walker equations of order $h$ and reads off the
coefficient of $X_{t+h-h} = X_t$ in the linear regression of $X_{t+h}$
on $X_{t+1}, \dots, X_{t+h}$. Most software (Python
\texttt{statsmodels}, R \texttt{stats}) returns the PACF in $O(h)$
operations per lag.
\end{rembox}

% =========================================================
\section{ACF / PACF signatures of standard processes}
\label{sec:acf-signatures}

\begin{center}
\begin{tabular}{lll}
\toprule
\textbf{Process} & \textbf{$\rho(h)$ behaviour} & \textbf{$\alpha(h)$ behaviour} \\
\midrule
$\WN$           & $\rho(h) = 0$ for $h \neq 0$         & $\alpha(h) = 0$ for $h \geq 1$ \\
$\AR{p}$        & decays (exp.\ / sinusoidal)          & cuts off at $h = p$ \\
$\MAm{q}$       & cuts off at $h = q$                  & decays \\
$\ARMA{p}{q}$   & decays after $h > q$                 & decays after $h > p$ \\
\bottomrule
\end{tabular}
\end{center}

\medskip
\noindent\textbf{Box--Jenkins identification.} Plot the sample ACF and
PACF, then use the table above to read off candidate orders $(p, q)$.

\subsection{AR(1) and MA(1) — closed-form ACF}

\begin{exbox}[AR(1) ACF]
\label{ex:ar1-acf}
For $X_t = \phi X_{t-1} + \varepsilon_t$ with
$\varepsilon_t \sim \WN(0, \sigma^2)$ and $\abs\phi < 1$:
\[
  \rho(h) = \phi^{\abs h}, \quad h \in \Z,
  \qquad
  \alpha(1) = \phi, \quad \alpha(h) = 0 \text{ for } h \geq 2.
\]
The ACF decays exponentially; the PACF cuts off at lag 1.
\end{exbox}

\begin{exbox}[MA(1) ACF]
\label{ex:ma1-acf}
For $X_t = \varepsilon_t + \theta \varepsilon_{t-1}$ with
$\varepsilon_t \sim \WN(0, \sigma^2)$:
\[
  \rho(0) = 1, \qquad
  \rho(1) = \frac{\theta}{1 + \theta^2}, \qquad
  \rho(h) = 0 \text{ for } h \geq 2.
\]
The ACF cuts off at lag 1; the PACF decays.
\end{exbox}

\begin{proof}[Proof of $\rho(1) = \theta / (1 + \theta^2)$ in MA(1)]
$\Var[X_t] = \E[(\varepsilon_t + \theta\varepsilon_{t-1})^2]
= (1 + \theta^2)\sigma^2$. Then
$\Cov(X_t, X_{t+1}) = \E[(\varepsilon_t + \theta\varepsilon_{t-1})
(\varepsilon_{t+1} + \theta\varepsilon_t)] = \theta\,\sigma^2$ since
$\varepsilon_t$'s are uncorrelated. Hence $\rho(1) =
\theta\sigma^2 / [(1 + \theta^2)\sigma^2] = \theta / (1 + \theta^2)$.
\end{proof}

% =========================================================
\section{Wold decomposition}
\label{sec:wold}

\begin{thbox}[Wold's theorem, 1938]
\label{th:wold}
Every \emph{purely non-deterministic} weakly stationary process
$\{X_t\}$ admits a unique representation
\[
  X_t \;=\; \mu \;+\; \sum_{j = 0}^{\infty} \psi_j\, \varepsilon_{t - j},
  \qquad
  \psi_0 = 1, \qquad \sum_{j = 0}^{\infty} \psi_j^2 < \infty,
\]
where $(\varepsilon_t)$ is a white-noise process (the
\emph{innovations}) of variance
$\sigma^2 = \Var[\varepsilon_t]$.
\end{thbox}

\begin{rembox}[Why it matters]
\begin{itemize}
  \item \textbf{Universality:} every stationary process can be viewed
        as an \emph{infinite MA process}.
  \item \textbf{Justification of ARMA:} ARMA models are
        finite-dimensional approximations of the Wold representation
        (with a rational $\psi(z)$).
  \item \textbf{Forecasting:} optimal linear predictors and their
        prediction errors are obtained from the Wold coefficients.
\end{itemize}
\end{rembox}

% =========================================================
\section{White-noise tests}
\label{sec:wn-tests}

After fitting a decomposition or an ARMA model, the residuals
$\widehat\varepsilon_t$ should look like white noise. Two classical
tests check this jointly across several lags.

\begin{defbox}[Box--Pierce and Ljung--Box statistics]
\label{def:lb}
For sample ACF $\widehat\rho(h)$ of the residuals up to lag $m$:
\[
  Q_{\text{BP}} = T \sum_{h = 1}^{m} \widehat\rho(h)^2,
  \qquad
  Q_{\text{LB}} = T(T + 2) \sum_{h = 1}^{m} \frac{\widehat\rho(h)^2}{T - h}.
\]
\end{defbox}

\begin{thbox}[Asymptotic distribution under $H_0$]
\label{th:lb}
Under $H_0 :$ \emph{the residuals are white noise}, both statistics
are asymptotically $\chi^2_{m - p - q}$, where $p + q$ is the number
of fitted ARMA parameters (zero if no parametric model has been
fitted).
\end{thbox}

\paragraph{Decision rule.} Reject $H_0$ at level $\alpha$ if
$Q > \chi^2_{m - p - q,\, 1 - \alpha}$, equivalently $p$-value $< \alpha$.

\paragraph{Choosing $m$.} A few rules of thumb:
\begin{itemize}
  \item For yearly data: $m \in \{10, 20\}$.
  \item For monthly data: $m \in \{12, 24, 36\}$ (covers seasonality).
  \item For high-frequency data: $m$ proportional to expected
        autocorrelation length.
\end{itemize}

\begin{warnbox}
The Ljung--Box test is preferred over Box--Pierce in finite samples:
its small-sample distribution is closer to the $\chi^2$ asymptotic
limit, especially when $m / T$ is not negligible.
\end{warnbox}

% =========================================================
\section*{Summary}
\addcontentsline{toc}{section}{Summary}

\noindent\textbf{Key takeaways.}
\begin{itemize}
  \item \emph{Strict stationarity} = invariance of all finite-dimensional
        distributions; \emph{weak stationarity} = only first two moments.
        Strict $\Rightarrow$ weak (under $L^2$); converse holds for
        Gaussian processes.
  \item The \emph{autocovariance} $\gamma(h)$ describes the dependence
        structure; the \emph{autocorrelation} $\rho(h)$ is its
        normalised version.
  \item ACF and PACF have characteristic patterns that allow
        \emph{Box--Jenkins identification} of AR(p) / MA(q) / ARMA(p,q).
  \item \emph{Wold decomposition}: every stationary process is an
        infinite moving average of innovations.
  \item \emph{Ljung--Box test} validates that residuals look like
        white noise.
\end{itemize}
